有效场论中的 Wilson 系数必须由高能理论的自由度和耦合决定，而不能仅凭低能算符的维数猜测。对任何在给定背景下只以二次型出现的重场，树级匹配首先是一个分块算符消元问题。把轻、重自由度组成列向量，并把二次核写成
\begin{equation*}
 \mathcal K(p)=
 \begin{pmatrix}
  K_{LL}(p)&K_{LH}(p)\\
  K_{HL}(p)&K_{HH}(p)
 \end{pmatrix}.
\end{equation*}
若过程能区内 $K_{HH}(p)$ 可逆，则消去重块以后，轻场的精确二次核为 Schur 补
\begin{equation}
 \boxed{
 K_{\rm eff}(p)
 =K_{LL}(p)-K_{LH}(p)K_{HH}^{-1}(p)K_{HL}(p).}
 \label{eq:sm-heavy-block-schur-matching-dp85}
\end{equation}
\eqnrefs{eq:sm-heavy-block-schur-matching-dp85} 保留 $K_{HH}^{-1}(p)$ 的完整动量依赖，是局域 EFT 展开之前的精确二次型母式。

设重块的最小质量能隙为 $M_{H,\min}^{(E)}$，过程特征能量为 $Q_E$。只有在
\begin{equation}
 \boxed{
 \epsilon_H\equiv\frac{Q_E}{M_{H,\min}^{(E)}}\ll1}
 \label{eq:sm-heavy-field-local-control-dp85}
\end{equation}
时，才可把重传播子按外部动量展开成局域算符级数。最低阶保留 $K_{HH}^{-1}(0)$，更高阶项依次带额外的 $Q_E/M_H^{(E)}$ 幂；这正是从一个确定 UV 理论得到低能 Wilson 系数的树级匹配程序。

\begin{exbox}[重右手中微子与 type-I seesaw]
在标准模型轻场之外加入规范单态右手中微子 $N_R$。沿用全书 Higgs--Yukawa 的 SI 归一化，其新增拉氏量为
\begin{equation*}
 \Delta\Lag_N
 =-\sqrt{\hbar c}\,\overline{L_L}Y_\nu\widetilde\Phi N_R
 -\frac{\hbar c}{2}\overline{N_R^{\,c}}\,\widehat M_R N_R
 +\mathrm{H.c.},
\end{equation*}
其中 $Y_\nu$ 是无量纲 Yukawa 矩阵，$\widehat M_R$ 具有逆长度量纲，并定义质量能量矩阵
\begin{equation*}
 M_R^{(E)}\equiv\hbar c\,\widehat M_R.
\end{equation*}
令
\begin{equation*}
 M_{R,\min}^{(E)}\equiv\sigma_{\min}(M_R^{(E)}),
 \qquad
 \epsilon_R\equiv\frac{Q_E}{M_{R,\min}^{(E)}}\ll1.
\end{equation*}
\eqnrefs{eq:sm-heavy-block-schur-matching-dp85} 中重 Majorana 传播子的轻--轻手征翻转部分在最低阶给出 $(M_R^{(E)})^{-1}$。代入两个 Yukawa 顶角，Weinberg 算符的能量归一化系数为
\begin{equation}
 \boxed{
 \frac{C^{(5)}}{\Lambda_E}
 =Y_\nu(M_R^{(E)})^{-1}Y_\nu^T+\delta C_5,
 \qquad
 \frac{\|\delta C_5\|}
 {\|Y_\nu\|^2\,\|(M_R^{(E)})^{-1}\|}
 =O(\epsilon_R^2).}
 \label{eq:typeI-weinberg-matching-dp86}
\end{equation}
Higgs 场取得真空期望值以后，$m_Dc^2=v_EY_\nu/\sqrt2$，轻中微子质量矩阵变成
\begin{equation}
 \boxed{
 M_\nu c^2
 =-\frac{v_E^2}{2}
 Y_\nu(M_R^{(E)})^{-1}Y_\nu^T
 +O\!\left(\frac{\|m_Dc^2\|^4}{[M_{R,\min}^{(E)}]^3}\right).}
 \label{eq:typeI-seesaw-mass-dp86}
\end{equation}
不同 UV 理论可以匹配到同一个 Weinberg 系数；type-I seesaw 对应其中由重 Majorana 单态产生的一组 Wilson 系数。
\end{exbox}
